\documentclass[12pt]{article}
\usepackage{amsmath, amssymb}
\usepackage[margin=1in]{geometry}
\setlength{\parskip}{6pt}
\title{QUBO Formulation: Hospital Staff Rescheduling Problem}
\date{}
\begin{document}
\maketitle

\subsection*{Hospital Staff Rescheduling Problem}

\paragraph{Problem Statement.}
Let $\mathcal{S}$ denote the set of available staff members and $\mathcal{J}$ denote the set of uncovered shifts. Let $\mathcal{A} \subseteq \mathcal{S} \times \mathcal{J}$ be the set of qualified staff-shift pairs, where $(s,j) \in \mathcal{A}$ indicates that staff member $s$ is qualified to cover shift $j$. For each qualified pair $(s,j) \in \mathcal{A}$, let $c_{s,j} \ge 0$ represent the cost of assigning staff $s$ to shift $j$, which may encode fatigue, travel distance, overtime, or preference penalties. The objective is to select a set of assignments such that every shift in $\mathcal{J}$ is covered by exactly one staff member, each staff member covers at most one shift, and the total assignment cost is minimized.

\paragraph{Decision Variables.}
For each qualified pair $(s,j) \in \mathcal{A}$, introduce a binary decision variable $x_{s,j} \in \{0,1\}$. The variable $x_{s,j}$ takes the value $1$ if staff member $s$ is assigned to shift $j$, and $0$ otherwise. The qualification constraint is enforced by restricting the variable domain to $\mathcal{A}$; variables corresponding to unqualified pairs are excluded from the formulation.

\paragraph{QUBO Derivation.}
We construct a Quadratic Unconstrained Binary Optimization (QUBO) formulation by converting the constrained minimization problem into an unconstrained energy minimization problem $H(\mathbf{x})$.

\textbf{Step 1: Objective Function.}
The original objective is to minimize the total assignment cost:
\begin{align}
    \min \sum_{(s,j) \in \mathcal{A}} c_{s,j} x_{s,j}.
\end{align}
This term is retained directly in the QUBO Hamiltonian.

\textbf{Step 2: Constraints.}
The problem imposes two types of constraints:
\begin{enumerate}
    \item \textit{Shift Coverage:} Each shift $j \in \mathcal{J}$ must be assigned exactly one staff member. Let $\mathcal{S}_j = \{s \in \mathcal{S} : (s,j) \in \mathcal{A}\}$ be the set of qualified staff for shift $j$. The constraint is:
    \begin{align}
        \sum_{s \in \mathcal{S}_j} x_{s,j} = 1 \quad \forall j \in \mathcal{J}.
    \end{align}
    \item \textit{Staff Limit:} Each staff member $s \in \mathcal{S}$ can cover at most one shift. Let $\mathcal{J}_s = \{j \in \mathcal{J} : (s,j) \in \mathcal{A}\}$ be the set of shifts qualified for staff $s$. The constraint is:
    \begin{align}
        \sum_{j \in \mathcal{J}_s} x_{s,j} \le 1 \quad \forall s \in \mathcal{S}.
    \end{align}
\end{enumerate}

\textbf{Step 3: Penalty Terms.}
We convert constraints to quadratic penalty terms using a penalty weight $\lambda > 0$.

\textit{Shift Coverage Penalty:}
For each shift $j$, we enforce equality using the squared penalty:
\begin{align}
    P_{\text{shift},j} = \lambda \left( \sum_{s \in \mathcal{S}_j} x_{s,j} - 1 \right)^2.
\end{align}
Expanding using the idempotent property $x_{s,j}^2 = x_{s,j}$ for binary variables:
\begin{align}
    P_{\text{shift},j} &= \lambda \left( \sum_{s \in \mathcal{S}_j} x_{s,j}^2 - 2 \sum_{s \in \mathcal{S}_j} x_{s,j} + 1 \right) \\
    &= \lambda \left( -\sum_{s \in \mathcal{S}_j} x_{s,j} + 2 \sum_{\substack{s,s' \in \mathcal{S}_j \\ s < s'}} x_{s,j} x_{s',j} + 1 \right).
\end{align}
Summing over all shifts gives the total shift penalty $H_{\text{shift}} = \sum_{j \in \mathcal{J}} P_{\text{shift},j}$.

\textit{Staff Limit Penalty:}
For each staff member $s$, we enforce the inequality $\sum_{j \in \mathcal{J}_s} x_{s,j} \le 1$ using the squared penalty $\lambda \left( \sum_{j \in \mathcal{J}_s} x_{s,j} \right)^2$, which is zero for sums $0$ or $1$ and positive for sums greater than $1$:
\begin{align}
    P_{\text{staff},s} = \lambda \left( \sum_{j \in \mathcal{J}_s} x_{s,j} \right)^2.
\end{align}
Expanding:
\begin{align}
    P_{\text{staff},s} &= \lambda \left( \sum_{j \in \mathcal{J}_s} x_{s,j}^2 + 2 \sum_{\substack{j,j' \in \mathcal{J}_s \\ j < j'}} x_{s,j} x_{s,j'} \right) \\
    &= \lambda \left( \sum_{j \in \mathcal{J}_s} x_{s,j} + 2 \sum_{\substack{j,j' \in \mathcal{J}_s \\ j < j'}} x_{s,j} x_{s,j'} \right).
\end{align}
Summing over all staff gives the total staff penalty $H_{\text{staff}} = \sum_{s \in \mathcal{S}} P_{\text{staff},s}$.

\textbf{Step 4: Hamiltonian.}
The total QUBO Hamiltonian combines the objective and penalty terms. We subtract a constant offset to remove the constant terms arising from the shift coverage expansion:
\begin{align}
    H(\mathbf{x}) = \sum_{(s,j) \in \mathcal{A}} c_{s,j} x_{s,j} + H_{\text{shift}} + H_{\text{staff}} - \lambda |\mathcal{J}|.
\end{align}

\textbf{Step 5: Q Matrix and Offset.}
Collecting coefficients for linear and quadratic terms yields the QUBO parameters. The Hamiltonian can be written as $H(\mathbf{x}) = \sum_{(s,j) \in \mathcal{A}} Q_{(s,j),(s,j)} x_{s,j} + \sum_{(s,j) \neq (s',j') \in \mathcal{A}^2} Q_{(s,j),(s',j')} x_{s,j} x_{s',j'} + c$. The entries are:
\begin{align}
    Q_{(s,j),(s,j)} &= c_{s,j} - \lambda + \lambda = c_{s,j}, \\
    Q_{(s,j),(s',j)} &= 2\lambda \quad \text{for } s \neq s', \\
    Q_{(s,j),(s,j')} &= 2\lambda \quad \text{for } j \neq j', \\
    Q_{(s,j),(s',j')} &= 0 \quad \text{otherwise}, \\
    c &= 0.
\end{align}
Here, $Q_{(s,j),(s',j)}$ corresponds to pairs sharing the same shift, and $Q_{(s,j),(s,j')}$ corresponds to pairs sharing the same staff member.

\paragraph{Penalty Weight Justification.}
The penalty weight $\lambda$ must be chosen sufficiently large to ensure that any violation of the constraints results in an energy increase that outweighs any potential reduction in the objective cost. Specifically, let $C_{\max} = \max_{(s,j) \in \mathcal{A}} c_{s,j}$. Choosing $\lambda > C_{\max}$ guarantees that the penalty for violating a constraint exceeds the maximum possible gain from the objective function. For the shift coverage constraint, a violation allows at most saving $C_{\max}$ in cost (by omitting an assignment), while incurring a penalty of at least $\lambda$. For the staff limit constraint, a violation incurs a penalty of at least $\lambda$ with no cost reduction. Thus, $\lambda > C_{\max}$ ensures that infeasible solutions have strictly higher energy than the best feasible solution.

\paragraph{Correctness Argument.}
Minimizing $H(\mathbf{x})$ over $\{0,1\}^{|\mathcal{A}|}$ recovers the optimal assignment schedule. For any feasible solution, the penalty terms vanish (after offset subtraction), so $H(\mathbf{x})$ equals the total assignment cost. For any infeasible solution, at least one penalty term is positive, contributing at least $\lambda$ to the energy. Since $\lambda > C_{\max}$, the energy of any infeasible solution exceeds the cost of any feasible solution (assuming costs are non-negative and at least one feasible assignment exists). Therefore, the global minimum of $H(\mathbf{x})$ must be a feasible solution with minimum total cost, corresponding to the optimal staff rescheduling.

\paragraph{Illustrative Example.}
Consider a small instance with $|\mathcal{S}| = 2$ staff members ($s \in \{0,1\}$) and $|\mathcal{J}| = 2$ shifts ($j \in \{0,1\}$), where all staff are qualified for all shifts, so $\mathcal{A} = \mathcal{S} \times \mathcal{J}$. Let the cost matrix be:
\begin{align}
    c_{0,0} = 10, \quad c_{0,1} = 20, \quad c_{1,0} = 15, \quad c_{1,1} = 5.
\end{align}
The decision variables are $x_{0,0}, x_{0,1}, x_{1,0}, x_{1,1}$. The QUBO Hamiltonian is:
\begin{align}
    H(\mathbf{x}) &= 10 x_{0,0} + 20 x_{0,1} + 15 x_{1,0} + 5 x_{1,1} \nonumber \\
    &\quad + \lambda (x_{0,0} + x_{1,0} - 1)^2 + \lambda (x_{0,1} + x_{1,1} - 1)^2 \nonumber \\
    &\quad + \lambda (x_{0,0} + x_{0,1})^2 + \lambda (x_{1,0} + x_{1,1})^2 - 2\lambda.
\end{align}
Expanding and collecting terms yields the Q matrix entries consistent with the general derivation. For $\lambda > 20$, the optimal solution minimizes the cost subject to constraints. The feasible assignments are $(x_{0,0}=1, x_{1,1}=1)$ with cost $15$, and $(x_{0,1}=1, x_{1,0}=1)$ with cost $35$. The global minimum is achieved at:
\begin{align}
    x_{0,0} = 1, \quad x_{1,1} = 1, \quad x_{0,1} = 0, \quad x_{1,0} = 0,
\end{align}
with minimum energy $H(\mathbf{x}^*) = 15$. This corresponds to assigning staff $0$ to shift $0$ and staff $1$ to shift $1$, yielding the minimum total cost of $15$.
\end{document}